\documentclass[11pt]{amsart}
\usepackage[T1]{fontenc}
\usepackage{iftex}
\ifPDFTeX\usepackage[utf8]{inputenc}\fi
\usepackage{lmodern}
\usepackage[margin=1.1in]{geometry}
\usepackage{amssymb,mathtools}
\usepackage{booktabs,array}
\usepackage{enumitem}
\usepackage{url}
\usepackage[hidelinks]{hyperref}
\hypersetup{pdftitle={Supplement: monostatic polyhedra},pdfauthor={Claudio Mattioni}}

\setlist{itemsep=2pt,topsep=4pt}
\renewcommand{\thesection}{S\arabic{section}}
\newtheorem*{lemma}{Lemma}
\newcommand{\conv}{\operatorname{conv}}
\newcommand{\vol}{\operatorname{vol}}
\urlstyle{tt}

\title{Supplement to ``Monostatic polyhedra: no mono-stable pentahedron, and \texorpdfstring{$9\le V_U\le14$}{9 <= V\_U <= 14}''}
\author{Claudio Mattioni}
\date{November 2026}

\begin{document}
\maketitle

This supplement collects experiments and side results that are not needed for Theorems A--C. Unless stated
otherwise they are numerical observations, not proofs. Every polyhedron that is called \emph{certified}
below was checked in exact rational arithmetic by \path{verify/exact_verify.py}. Section, lemma and
proposition numbers without the prefix S refer to the main paper.

\section{The vertex lower bound at \texorpdfstring{$V=9$}{V = 9}}

Proving $V_U\ge10$ requires excluding $V=9$. There are 50 simplicial 2-spheres with 9 vertices; the list
is complete by the flip-closure check (\path{lb_joint/flip_closure.py}, 667 flips). There are 9,915,284
(sphere, root, tree) cases.

\emph{Stage 1} (\path{lb_joint/run_V9_stage1.log}). The single-cone certificate of Proposition~4.2 excludes
9,163,771 cases (92.4\%). The certificate file is 1.46\,GB and is not part of the repository.

\emph{Stage 2.} For the remaining 751,513 trees three further sound ingredients were tried. All of them are
implemented, with their soundness arguments, in \path{lb_joint/verify_joint.py}.
\begin{enumerate}
\item \emph{Coupled cones} (record kind \texttt{certG}; Proposition~4.3).
\item \emph{Hybrid certificates} (record kind \texttt{certH}). The triangles avoiding $v$ are split into
  ``good'' and ``bad'' ones. A good $T$ needs a positive definite block of Proposition~4.2. A bad $T$
  needs, for every tetrahedron $T'$ of the cone decomposition from a second apex $q$, a positive definite
  block that combines both directions $\mu$ and $\mu'$. This is about 4 times faster than \texttt{certG},
  because no coupling matrix is needed.
\item \emph{Steepest-neighbour inequalities} (S1). Assign to each vertex $i\ne v$ its steepest shadowing
  neighbour $p(i)\in\arg\max_{k\sim i}r_k\cdot r_i$. This is again an acyclic shadowing tree, and it
  satisfies $r_i\cdot r_k\le r_i\cdot r_{p(i)}$ for all neighbours $k$. Every configuration has such a tree,
  so these inequalities may be added for each enumerated tree.
\end{enumerate}

\emph{Obstruction} (\path{lb_joint/v9_obstruction.py}, \texttt{.log}, \texttt{.json}). We sampled 3,000 of
the stage-2 trees uniformly (seed 2026):
\begin{center}
\begin{tabular}{lr}
\toprule
outcome & trees\\
\midrule
excluded by a hybrid certificate & 2,911\\
excluded by all nine cone decompositions + S1 after the earlier searches failed & 22\\
left uncertified by the tested searches & 67\\
\bottomrule
\end{tabular}
\end{center}
For the 67 survivors, the search with all nine cone decompositions and the S1 inequalities did not produce
an exact certificate. The script records the numerical SDP objective. It accepts this last certificate
type only when that objective exceeds $10^{-7}$ and integer rounding passes the exact check. It does not
construct or verify an exact feasible point of the relaxation for the failures. These outcomes therefore
do not establish feasibility of the relaxation or non-existence of certificates of the tested types.

Extrapolating this sample suggests that about 17,000 trees (0.17\% of all cases) would remain uncertified
by the same procedure; this is not an exhaustive count. All survivors in the sample have maximal tree
depth 6 or 7. A typical one is sphere 0, root 0, with the path
$0\leftarrow7\leftarrow6\leftarrow5\leftarrow4\leftarrow3\leftarrow2\leftarrow1$ and the leaf
$0\leftarrow8$. The 294 sampled trees of depth 8 were all excluded.

The relaxation never uses the actual volume weights of the tetrahedra (only their nonnegativity) or the
convexity of the configuration beyond the shadowing inequalities. These are the natural places to look
for stronger valid constraints.

\section{Thin plates in more detail}

\subsection{Flat pyramids: 16 vertices}
Take the outline polygon $B$ in a plane and one apex lifted slightly above a point $A'\in B$. The centre of
mass of the pyramid projects to $(3g+A')/4$, where $g$ is the centroid of $B$. Put $A'$ at the bottom of a
double tight spiral and optimise the spiral angles. This works for outlines with 15 vertices (arm splits
7+7, 8+6, 9+5, 10+4), giving certified mono-unstable polyhedra with 16 vertices
(\path{verify/certificates/seed_V16*}). For 14-gons no feasible candidate was found in the tested arm
splits. The previous record, 18 vertices~\cite{DKLRV20}, is a flat pyramid over Conway and Guy's 17-gon;
certified versions are in \path{verify/certificates/seed_V18*}.

\subsection{Plates with heights on the outline}
The construction of Theorem~C (\S5.2) needs no apex. CMA-ES on the planar surrogate (the thin-plate limit
of \S5.1) and then in full 3D finds 14-vertex examples for arm splits 7+6, 8+5 and 9+4. A version with a
larger safety margin is \path{search_vertices/certified/V14_spiral76_d7.json}. A certified 15-vertex body
is \path{search_vertices/certified/V15_insert_d7.json}.

\subsection{Thirteen vertices: negative evidence for thin plates}
\begin{itemize}
\item \emph{Search.} 328 optimisation runs with about $3\cdot10^7$ evaluations covered all arm splits and
  many height patterns. They ended at clean minimax points with all margins equal and negative: the best
  angular margin was $-7.9\cdot10^{-3}$, against $+1.07\cdot10^{-2}$ at 14 vertices.
\item \emph{Flat limit.} In the reported three-dimensional optimisation runs, the thickness tends towards
  zero. This is evidence favouring thin plates within the tested searches, not a proof of local
  optimality.
\item \emph{Frustum thickness model.} \path{explore/ideal_frustum.py} searches bodies of the form
  $\conv((K\times\{0\})\cup(E\times\{1\}))$, where $E$ is the portion of a spiral outline $K$ below a
  chosen level. Their thickness has a plateau on $E$ and a concave extension over the rest of $K$. The
  script optimises the level and the spiral angles numerically. It does not prove that this thickness is
  optimal among all admissible functions, nor that the numerical minima are global. In units with
  $|T-O|=1$, the best reported centroid heights are $+0.0285$ for an 11-gon outline and $-0.0034$ for a
  13-gon outline. The construction can add vertices beyond those of the outline, so these observations do
  not give a lower bound on the vertex count of a mono-unstable polyhedron.
\end{itemize}
None of this is a proof, and genuinely three-dimensional (non-plate) designs were explored only lightly.

\section{Polarity between the two record bodies}

\begin{lemma}[polarity]
Let $0\in\operatorname{int}P$, $P=\{x:\ell_i\cdot x\le1\}$, and let $P^\circ=\conv\{\ell_i\}$ be the polar
body. Face $i$ of $P$ is stable with respect to the origin iff $\ell_i$ is an unstable vertex of $P^\circ$
with respect to the origin. Both conditions say $\ell_j\cdot\ell_i<|\ell_i|^2$ for all $j\ne i$
(Lemma~2.2).
\end{lemma}

So stable faces of $P$ correspond to unstable vertices of $P^\circ$, but with respect to the centre of
mass of $P$, which in general is not the centre of mass of $P^\circ$. This is why the face problem and the
vertex problem are different problems.

Reshetov's 14-face mono-stable polyhedron is a long thin spindle whose 14 faces are all ``lateral'' and
slightly tilted. Its polar about the centre of mass is a thin plate with 14 vertices, mono-unstable with
respect to the origin. The centre of mass of that plate misses the origin by only 0.7\% of the diameter,
and a local optimisation turns it into a genuine mono-unstable polyhedron
(\path{search_vertices/certified/V14_resh_d6.json}). In the limits ``infinitely long spindle'' and
``infinitely thin plate'' both problems become planar problems about a double spiral. They differ only in
which point must lie below the spiral centre: the centroid of the plate, or the Santal\'o point of the
plate.

\section{Faces: Reshetov's bodies and improved margins}

Reshetov's four published bodies, with 14--17 faces and integer plane data~\cite{Res14}, were re-verified
exactly (\path{search_faces/certified/reshetov1*.json}). We measure the safety margin as
in~\cite{Res14}: the distance of the centroid projections from the face boundaries, divided by
$\vol^{1/3}$. On this measure the published bodies have margins $1.4\cdot10^{-4}$ (14 faces) and
$6.4\cdot10^{-4}$ (15 faces); his reported best values are $1.2\cdot10^{-3}$ and $4.2\cdot10^{-3}$.

Local optimisation reaches $2.2\cdot10^{-3}$ (14 faces) and $5.5\cdot10^{-3}$ (15 faces) in floating
point. The bodies \path{search_faces/certified/improved14.json} and \path{improved15.json} are rounded to
4-digit integer plane data and are exactly certified mono-stable, with margins $1.2\cdot10^{-3}$ and
$4.6\cdot10^{-3}$. About 100 optimisation runs aimed at 13 faces failed. This is a small effort compared
with Reshetov's $10^{10}$ evaluations and proves nothing.

\section{Simplices}

For $d$-simplices, the conditions of \S3.1 involve only the Gram matrix of the scaled facet normals, which
has zero row sums. ``Is there a monostatic $d$-simplex with a given rolling pattern'' then reduces to one
small semidefinite feasibility problem. All 719 rooted trees on 10 nodes are infeasible numerically,
consistent with the absence of a monostatic 9-simplex. This case was left open by Dawson, Finbow and Mak.
It is not new: it is equivalent to \cite[Cor.~6]{PRDB23}, which states that mono-unstable 0-skeleta need 11
vertices in every dimension.

\section{Exploration code}

\begin{center}\small
\begin{tabular}{>{\raggedright}p{0.3\textwidth}p{0.62\textwidth}}
\toprule
directory & content\\
\midrule
\path{lowerbound/} & replication of the SDP method of~\cite{BDPR24}, and where it stops at $V=7$\\
\path{lb_joint/} & certificate generators (\path{certify_joint.py}, \path{general_cert.py},
  \path{hybrid_cert.py}, \path{stage2_hybrid.py}), the trusted verifier and the $V=9$ experiments\\
\path{penta/} & \path{search.py}, \path{propose.py}, \path{compact.py} (untrusted; they produce the cover);
  \path{sanity.py} (Monte-Carlo test)\\
\path{search_vertices/}, \path{search_faces/}, \path{explore/} & optimisation searches and prototypes\\
\path{audit/} & independent AI audits of Theorem~A (\path{REPORT.md}) and Theorem~B
  (\path{theoremB/REPORT.md}), with separate checkers and logs\\
\bottomrule
\end{tabular}
\end{center}

\begin{thebibliography}{DKLRV20}
\bibitem[BDPR24]{BDPR24} S.~Boz\'oki, G.~Domokos, D.~Papp, K.~Reg\H{o}s, \emph{The smallest mono-unstable,
homogeneous convex polyhedron has at least 7 vertices}, arXiv:2401.17906 (2024).
\bibitem[DKLRV20]{DKLRV20} G.~Domokos, F.~Kov\'acs, Z.~L\'angi, K.~Reg\H{o}s, P.~T.~Varga, \emph{Balancing
polyhedra}, Ars Math.\ Contemp.\ 19 (2020) 95--124.
\bibitem[PRDB23]{PRDB23} D.~Papp, K.~Reg\H{o}s, G.~Domokos, S.~Boz\'oki, \emph{The smallest mono-unstable
convex polyhedron with point masses has 8 faces and 11 vertices}, Eur.\ J.\ Oper.\ Res.\ 310 (2023)
511--517.
\bibitem[Res14]{Res14} A.~Reshetov, \emph{A unistable polyhedron with 14 faces}, Int.\ J.\ Comput.\ Geom.\
Appl.\ 24 (2014) 39--59.
\end{thebibliography}

\end{document}
